Let \(z_0\) be a zero of \(M\) of multiplicity \(\ell\). By the definition of multiplicity,
\[
M^{(j)}(z_0)=0\quad(0\le j<\ell),
\qquad
M^{(\ell)}(z_0)\ne0.
\]
The moment generating function \(M\) is entire by its declared definition. On \(\mathcal P_{\mathrm{NG},n}\), the sub-Gaussian bound on \(\eta\), the range restrictions \(\lVert g_0\rVert_\infty\le C_g\) and \(\bar g_n\in[-C_g,C_g]\), and the conditional sub-Gaussian bound on \(\xi\) ensure that every polynomial--exponential moment below is finite. Thus differentiation under the relevant expectations is justified by dominated convergence on each compact subset of \(\mathbb C\).

Fix \(d\in\mathbb R\), and define \(L_d(z)=E_P[e^{z(\eta+d)}]\). Since \(d\) is deterministic,
\[
L_d(z)=e^{zd}M(z).
\]
The derivative representation of an entire moment generating function and the Leibniz rule therefore give
\[
\begin{aligned}
E_P[J_{z_0,\ell}(\eta+d)]
&=E_P[(\eta+d)^{\ell-1}e^{z_0(\eta+d)}]\\
&=L_d^{(\ell-1)}(z_0)\\
&=e^{z_0d}\sum_{j=0}^{\ell-1}{\ell-1\choose j}
d^{\ell-1-j}M^{(j)}(z_0)\\
&=0.
\end{aligned}
\]
The last equality follows from the multiplicity relations. It includes the boundary case \(\ell=1\), for which the sum consists only of \(M(z_0)=0\).

Now fix \(P\in\mathcal P_{\mathrm{NG},n}\). The range restrictions imply
\[
|D_n|=|g_0(X)-\bar g_n(X)|\le 2C_g
\quad P\text{-almost surely}.
\]
The random variable \(D_n\) is measurable with respect to \(X\), while \(\eta\perp X\) by \(\text{ass:independent-treatment-noise}\). Conditional on \(X\), apply the deterministic-shift identity with \(d=D_n\) to obtain
\[
E_P[J_{z_0,\ell}(Z_n)\mid X]
=E_P[J_{z_0,\ell}(\eta+D_n)\mid X]
=0
\quad P\text{-almost surely}.
\]

Independence of \(\eta\) and \(D_n\) also yields, for every \(z\in\mathbb C\),
\[
F_n(z)=E_P[e^{z(\eta+D_n)}]=M(z)H_n(z).
\]
Because \(Z_nJ_{z_0,\ell}(Z_n)=Z_n^\ell e^{z_0Z_n}\), differentiation followed by the Leibniz rule gives
\[
\begin{aligned}
E_P[Z_nJ_{z_0,\ell}(Z_n)]
&=F_n^{(\ell)}(z_0)\\
&=\sum_{j=0}^{\ell}{\ell\choose j}
M^{(j)}(z_0)H_n^{(\ell-j)}(z_0)\\
&=M^{(\ell)}(z_0)H_n(z_0).
\end{aligned}
\]
Every term with \(j<\ell\) vanishes, and the term \(j=\ell\) is the sole remaining term. Since \(M^{(\ell)}(z_0)\ne0\), the explicit relevance boundary condition is equivalently \(H_n(z_0)\ne0\); no further overlap or rank condition is used.

Finally, the residual definitions give
\[
Y=\theta_0Z_n+b_n(X)+\xi,
\qquad
b_n(X)=q_0(X)-\theta_0D_n.
\]
The class range restrictions make \(b_n\) bounded. Hence conditional annihilation implies
\[
E_P[b_n(X)J_{z_0,\ell}(Z_n)]
=E_P\!\left[b_n(X)E_P[J_{z_0,\ell}(Z_n)\mid X]\right]
=0.
\]
Moreover, \(J_{z_0,\ell}(Z_n)\) is measurable with respect to \((X,T)\), because \(Z_n=T-\bar g_n(X)\). By \(\text{ass:outcome-mean-independence}\),
\[
E_P[\xi J_{z_0,\ell}(Z_n)]
=E_P\!\left[J_{z_0,\ell}(Z_n)E_P[\xi\mid X,T]\right]
=0.
\]
Multiplying the displayed structural decomposition by \(J_{z_0,\ell}(Z_n)\) and taking expectations therefore yields
\[
E_P[YJ_{z_0,\ell}(Z_n)]
=\theta_0E_P[Z_nJ_{z_0,\ell}(Z_n)].
\]
Whenever \(E_P[Z_nJ_{z_0,\ell}(Z_n)]\ne0\), division is valid and proves
\[
\theta_0
=\frac{E_P[YJ_{z_0,\ell}(Z_n)]}
{E_P[Z_nJ_{z_0,\ell}(Z_n)]}.
\]
Thus the ratio is an identifying functional for the causal coefficient under the stated structural assumptions and the explicit nonzero-relevance boundary condition; the equality is derived rather than imposed by definition.